\begin{remark}[从递推系数还原 \texorpdfstring{$P_{16},P_{17}$}{P16,P17}]\label{rem:pom-moment-charpoly-16-17}
沿用定理 \ref{thm:pom-moment-resonance} 的递推格式与特征多项式约定
\(
P_q(\lambda)=\lambda^d-\sum_{i=1}^d c_i\lambda^{d-i}
\)。
由表 \ref{tab:fold_collision_moment_recursions_9_17} 的 \(q=16,17\) 递推系数直接代入可得
$$
\begin{aligned}
P_{16}(\lambda)={}&\lambda^{13}-2\lambda^{12}-1611\lambda^{11}-62312\lambda^{10}-2559407\lambda^9\\
&+24862788\lambda^8+585266591\lambda^7-62312\lambda^6+44606766\lambda^5\\
&-24862794\lambda^4-255692\lambda^3+124624\lambda^2-8\lambda+8,
\\
P_{17}(\lambda)={}&\lambda^{13}-2\lambda^{12}-2599\lambda^{11}-125872\lambda^{10}-6569850\lambda^9\\
&+96034590\lambda^8+2764163954\lambda^7+643026032\lambda^6+15022392733\lambda^5\\
&-769974566\lambda^4-15329386299\lambda^3-642908352\lambda^2-1347896340\lambda+673948170.
\end{aligned}
$$
这些显式特征多项式可直接用于根扫描/谱包结构分析；例如表 \ref{tab:fold_collision_kernel_rh_scan_q2_17} 的计算即以 \(P_q\) 为输入（脚本 \texttt{scripts/exp\_fold\_collision\_kernel\_rh\_scan\_q2\_17.py}）。
\end{remark}

